\documentclass[11pt,a4paper]{article}
\usepackage[T1]{fontenc}
\usepackage{lmodern}
\usepackage[english]{babel}
\usepackage{amsmath,amssymb,mathtools}
\usepackage[margin=2.2cm]{geometry}
\usepackage{setspace}
\usepackage[hidelinks]{hyperref}
\setstretch{1.25}
\frenchspacing
\emergencystretch=3em
\setlength{\parindent}{0pt}
\renewcommand{\le}{\leqslant}
\renewcommand{\ge}{\geqslant}
\DeclareMathOperator{\tr}{tr}
\DeclareMathOperator{\Tr}{Tr}
\DeclareMathOperator{\Pf}{Pf}
\newcounter{fact}
\newcommand{\fact}[2]{\par\bigskip\stepcounter{fact}\noindent\textbf{Fact \thefact\ (#1).}\ #2}
\newcommand{\used}[1]{\par\smallskip\noindent{\small\textit{Used in:} #1}\par}
\begin{document}
\begin{center}
{\LARGE\bfseries Linear algebra -- new problems\par}\vspace{0.6em}
{\large Supplement to the theoretical minimum\par}
\end{center}
\medskip
Facts needed for the problems in this folder that are not stated in the theoretical minimums of the sheets and cannot
be avoided by a solution with the methods of the sheets.
Problem numbers refer to \emph{Linear algebra - proposals.pdf}.

\fact{Finite fields}{For every prime power $q$ and every $n\ge1$ there is an irreducible polynomial $f$ of degree
$n$ over $\mathbb F_q$: the number of monic irreducible polynomials of degree $n$ equals
$\frac1n\sum_{d\mid n}\mu(d)q^{n/d}>0$. Then $\mathbb F_q[x]/(f)$ is a field with $q^n$ elements; via the companion
matrix of $f$ it is realised inside $M_n(\mathbb F_q)$ as a set of $q^n$ matrices in which every non-zero matrix is
invertible. (The sheets state this only for $q=p$ prime.)}
\used{Problem 18.}

\end{document}
